% =====================================================================
\section{How this is known}
Almost everything said above was obtained not in sick people. The closer a method of measurement is to the molecule, the further it is from the disease; the reliability of a conclusion is determined by the rung on which it was obtained (table~1).
\begin{table*}[t]\small\centering
\begin{tabular}{@{}>{\raggedright\arraybackslash}p{0.17\textwidth}>{\raggedright\arraybackslash}p{0.37\textwidth}>{\raggedright\arraybackslash}p{0.38\textwidth}@{}}
\toprule
\textsf{\textbf{Rung}} & \textsf{\textbf{What it permits}} & \textsf{\textbf{What has to be paid}}\\
\midrule
Purified proteins and cell lines & measuring reaction rates; hence the kinetics of peroxiredoxin & these are not neurons; the conditions are far from a living cell\\
\addlinespace
Neuronal culture & following cells for hours, adding substances at an exact dose & the cells are taken from embryos or newborns, they are immature and deprived of their usual surroundings\\
\addlinespace
Brain slice & recording the current of an individual cell in an intact local network & a slice lives for hours; long-range connections are cut\\
\addlinespace
Living animal & intervening and observing the consequences over days and months; knocking out individual genes & a mouse does not have schizophrenia; a link is reproduced, not the disease\\
\addlinespace
Human: post-mortem tissue & counting cells and proteins in the region of interest & a terminal state after decades of illness and treatment\\
\addlinespace
Human: in-vivo imaging & measuring in patients and over time & an average over a block of tissue of cubic centimetres\\
\bottomrule
\end{tabular}
\caption{Rungs of evidence: from the molecule to the human being.}
\end{table*}
\facts
\begin{itemize}
\item \emph{Antibody staining}: a tissue slice is treated with a colour-labelled antibody against the protein of interest, and cells or brightness are counted. A cell that has ceased producing the labelled protein looks as though it has vanished. ``Loss of PV cells'' can therefore mean either death or silence.\tU
\item \emph{Separation of proteins in a gel}: proteins are separated by size, labelled with an antibody, and the brightness of the band is measured. The result is relative, usually as a fraction of a control or of a maximum.\tU
\item \emph{Recording a cell's current}: a fine electrode on the membrane of a single cell in a slice; it yields exact numbers for a few cells.\tU
\item \emph{Radioactive-label imaging} (PET) shows, for example, the capacity to synthesise dopamine; \emph{magnetic resonance spectroscopy} shows the concentration of a substance, glutathione in particular, averaged over all cells of the volume.\tU
\item Groups in animal experiments are small: three to ten animals per point is commonplace.\tU
\end{itemize}
\begin{ineq}
For model building three rules follow from this. A number measured on one rung cannot be substituted on another without qualification. Relative units yield ratios and phase shifts, but not absolute values. An average over a volume of tissue says nothing about a rare cell type: PV cells make up a few per cent of the neurons.
\end{ineq}
